% !TEX root = ../../../tree_notes.tex
% =====================================================================
%  14.1  Tree knapsack and the amortised child-merge   sec:c27-knapsack
%  Exports: def:c27-knap-state  thm:c27-amortize  lst:c27-knap
%           fig:c27-mergetree  ex:c27-knap-run  prob:c27-knap-solved
%           lst:c27-855c-soln
%  THE owed proof: Ch.13 (ch10 / c10) previewed the O(nV) budgeted
%  child-merge twice and never proved it.  That proof is thm:c27-amortize
%  below, stated exactly as promised.  Anchor problem CF 855C (Helga
%  Hufflepuff's Cup), a three-class budgeted tree DP over j = #type-k
%  vaults with j <= x <= 10 -- a textbook V = x knapsack.  All merge
%  counts hand-computed on the running tree T from the opener.
% =====================================================================
\section{Tree Knapsack: an Index on the State, and the Merge Bound}\label{sec:c27-knapsack}

\begin{previouslybox}
\begin{enumerate}[nosep, leftmargin=1.6em]
  \item Chapter~13 folded a \emph{constant} tuple $\mathcal S(v)$ through one
        child at a time (\thmref{c10-subtree-template}). What made that
        $\bigO{n}$ rather than $\bigO{n^2}$? (each node did $\bigO{\deg(v)}$
        work; the children's tuples were combined independently, no shared
        budget)
  \item Chapter~13's distance-$k$ histogram (\secref{c10-distance}) stored an
        \emph{array} $H_v[\cdot]$ and merged children by a double loop. Why
        did it not need a full $\bigO{n^2}$ argument? (the array length was
        capped by $k$, a small constant -- the seed of today's bound)
  \item Chapter~12 computed $\size(v)$ for every node in one DFS
        (\coderef{c09-dfs3}). Read back the running tree's sizes
        $\size=[10,4,5,1,2,1,3,1,2,1]$. Which single quantity will bound our
        loop nests today? (the \emph{subtree size}, capped by the budget)
  \item State Ch.1's convention that made a leaf fall out of the fold with no
        \code{if}. (\term{$\height(\nil)=-1$}, $\size(\nil)=0$; today it
        becomes $\mathit{dp}[\ell][0]=0$ for a leaf)
\end{enumerate}
\end{previouslybox}

\begin{topicintro}
Chapter~13's recipe broke the moment the state stopped being a fixed-size
bag. Give each node an \emph{array} $\mathit{dp}[v][0\,.\,V]$ -- a value for
every amount of budget $j$ from $0$ to $V$ -- and children can no longer be
combined independently, because a unit of budget spent inside child $A$'s
subtree is a unit \emph{not} available to child $B$. The children must be
merged \emph{one at a time}, each merge a small knapsack convolution, and the
naive reading of the loops says the whole thing is quadratic (subtree sizes
multiply!). It is not. This section writes the merge template and proves, in
full, the bound Chapter~13 promised and never delivered:

\term{the amortised child-merge} (\thmref{c27-amortize}) -- bounding each
merge's two loop heads by the \emph{running} size makes the total
$\bigO{nV}$ across all nodes. CF~855C (\probref{c27-knap-solved}) is the
worked example: colour a tree with $m$ types, at most $x$ of them the
distinguished type~$k$, and count the colourings -- a knapsack whose budget
is ``number of type-$k$ vaults spent so far.''

\textbf{Prerequisites.} Chapter~13's child-fold template
\thmref{c10-subtree-template} and histogram merge (\secref{c10-distance});
Chapter~12's subtree sizes $\size(v)$ (\coderef{c09-dfs3}) -- used here as the
loop bound, not re-derived; \chref{c04} rooting by parent parameter.

\textbf{After this section you can}
\begin{itemize}[nosep]
  \item write the budgeted child-merge skeleton and keep every inner loop
        capped by the running size; \hfill\lv{2}
  \item state and prove \thmref{c27-amortize} -- Ch.13's outstanding debt; \hfill\lv{4}
  \item cold-write CF~855C's three-class cup-counting DP and trace it on an
        official sample; \hfill\lv{4}--\lv{5}
\end{itemize}
\end{topicintro}

\subsection{Intuition}

A scalar state composes with $\max$ or $+$ and the children stay independent:
$\mathcal S(v) = \mathcal F(\{\mathcal S(c)\})$ spends $\bigO{\deg(v)}$, so
$\sum_v \bigO{\deg(v)} = \bigO{n}$ (Ch.13). Put an index on the state and
independence dies. If $\mathit{dp}[v][j]$ means ``best value when the subtree
of $v$ consumes $j$ units of a global budget of $V$,'' then to close
$\mathit{dp}[v][\cdot]$ we must decide how to \emph{split} $v$'s budget
across its children, and every split $(j_1, j_2, \dots)$ with
$\sum_c j_c \le j$ is a candidate. That is a \term{knapsack}: we cannot feed
all children to $\mathcal F$ at once, so we fold them \emph{sequentially},
maintaining a running ``budget used so far at $v$'' array and convolving one
child in at a time.

The loop nest for one merge looks like two nested $[0..\text{size}]$ ranges,
which screams quadratic -- and if you write it naively (heads at $\text{sz}(v)$
and $\text{sz}(c)$) a chain input \emph{really is} $\Theta(n^2)$. The saving
grace is that a budget of $V$ can never be split past $V$: the reachable
budget at any node is $\min(\text{size so far}, V)$, so both loop heads are
$\min(\cdot, V)$. \thmref{c27-amortize} shows the total over every node and
every child of these capped products is only $\bigO{nV}$, because a merge step
is a pair of budget slots that is charged at most once (their lowest common
ancestor is the merging node), and the cap puts at most $V$ live slots per
node. Chapter~13 asserted this without proof; we prove it now.

\subsection{Definitions and Notation}

\begin{definition}{Budgeted subtree state and child merge}{c27-knap-state}
Fix a root and let $\children(v)$ be $v$'s children. A \term{budgeted subtree
state} is an array $\mathit{dp}[v][0\,.\,V]$ where $\mathit{dp}[v][j]$ is the
optimal value (or, for counting DPs, the number of configurations) of
$\subtree(v)$ using \emph{at most} $j$ units of budget, subject to any local
constraint the problem imposes. A leaf has $\mathit{dp}[\ell][0]$ fixed by its
own pick and higher $j$ weakly better. Children are combined by the
\term{merge}: folding a child $c$ into $v$ is the convolution
\[
  \mathit{dp}_{\text{new}}[v][t] \;=\; \max_{\,i+j\,=\,t}\Big(\mathit{dp}[v][i] + \mathit{dp}[c][j]\Big)
\]
(for maximisation; replace $\max\!\sum$ by $\sum\!\prod$ for counting), where
$i$ ranges over the budget already reachable at $v$ and $j$ over the budget
child $c$ can supply. The \term{answer} is read at the root, $\max_{j\le V}\mathit{dp}[r][j]$
(or its counting analogue). We require the merge to be defined on a leaf
(no children): $\mathit{dp}[v][\cdot]$ starts at $v$'s own contribution.
\end{definition}

\noindent\textbf{In plain words.} Each node remembers one number \emph{per
budget level}; a child is merged by trying every way to split the budget
between ``what $v$ already holds'' and ``what this child adds,'' keeping the
split at or below $V$. The order of children does not change the result --
convolution over $\le V$ is associative and commutative -- but merging one at
a time is what lets us cap the loops.

\subsection{The Merge Bound: Proving Ch.~13's Debt}

\begin{theorem}{The amortised child-merge costs $\bigO{nV}$}{c27-amortize}
Run the merge of \defref{c27-knap-state} over a rooted tree of $n$ nodes, and
at each node $v$ merge its children in any fixed order, letting $S_{v,c}$
denote the number of nodes already absorbed into $v$ just before child $c$
(so $S_{v,c} = 1 + \sum_{c' \text{ before } c} \size(c')$). If the two inner
loop heads are capped as
\[
  0 \le i \le \min(S_{v,c}, V), \qquad 0 \le j \le \min(\size(c), V), \qquad i + j \le V,
\]
then the total number of $(i,j)$ iterations over all nodes and all children is
$\bigO{nV}$.
\end{theorem}

\begin{proofbox}[\thmref{c27-amortize}: charge each iteration to a pair, each pair once, $\le V$ pairs per node]
An iteration is a tuple $(v, c, i, j)$ with $i + j \le V$. Charge it to the
\emph{ordered pair of nodes} $(x, y)$ defined as follows: $y$ is the $j$-th
node of $\subtree(c)$ in a fixed DFS order (with $y = \varnothing$ when $j = 0$),
and $x$ is the $i$-th node of $v$'s already-absorbed region
$\{v\} \cup \bigcup_{c' \text{ before } c} \subtree(c')$ in the same order
($x = \varnothing$ when $i = 0$). Because $i \le \min(S_{v,c}, V)$ and
$j \le \min(\size(c), V)$, both $x$ and $y$ are well-defined and lie inside
the capped first-$V$ window of their regions, so the map is injective: no two
tuples with different $(i,j)$ pick the same $(x,y)$.

\emph{Each pair is charged at most once.} Fix distinct $x, y$ and let
$a = \lca(x, y)$. A merge can charge $(x,y)$ only when $x$ and $y$ sit on
opposite sides of the merge: $y$ inside the child $c$ being folded and $x$
either equal to the merging node or inside an already-folded child. The only
node at which $x$ and $y$ are separated in exactly that way is their lowest
common ancestor $a$; at every ancestor strictly above $a$ both lie in the same
child subtree (so they are never split), and at every node below $a$ at least
one of them is not in the subtree at all. Hence at most one merge -- the one
at $a$, with the child $c$ being the child of $a$ that contains $y$ -- ever
charges $(x,y)$.

\emph{The cap bounds each node's charges by $V$.} For a fixed $y$, the tuples
that charge a pair ending in $y$ have $y$ as the $j$-th DFS node of $c$ for
$c$ on the path $a \to y$; there is a unique such $c$ per candidate ancestor
$a$, and $j$ is then determined. Since $i + j \le V$ forces $j \le V$, and the
$V$-capped window means $y$ is inside the first-$V$ nodes of $\subtree(c)$ only
while $a$ is within $V$ levels of the top of that child subtree, the number of
distinct merges at which $y$ is charged is at most $V$. Thus $y$ appears in at
most $V$ charged pairs; summing over the $n$ choices for $y$ gives at most
$nV$ charged pairs, i.e.\ $\bigO{nV}$ iterations. \qed
\end{proofbox}

\begin{insightbox}
The theorem is why Ch.~13's distance-$k$ DP ($V = k$) and today's cup DP
($V = x \le 10$) are both called ``linear'': when the budget $V$ is a small
constant, $\bigO{nV}$ is $\bigO{n}$. The cap \emph{is} the whole game. Write
the loop head as $\text{sz}(v)$ instead of $\min(\text{sz}(v), V)$ and the
chain input reverts to $\Theta(n^2)$ -- and it will still pass every hand
sample, because those samples are trees of four nodes. This is the number one
source of ``accepted locally, TLE on test 12'' in budgeted tree DP.
\end{insightbox}

\subsection{The Algorithm: The Budgeted Child-Merge}

\begin{lstlisting}[caption={The generic budgeted child-merge: same DFS as \coderef{c10-fold}, but the state is an array and children fold by capped convolution}, label={lst:c27-knap}]
// Budgeted subtree DP. dp[v][j] = value of subtree(v) using <= j budget.
// Merge children ONE AT A TIME. Cap BOTH loop heads by the running sizes so
// the whole DFS costs O(nV) -- thm:c27-amortize -- not O(n^2).
int V;                                   // per-node budget cap
vector<vector<int>> adj;
vector<int> sz;                          // subtree size (Ch.12 metric)
vector<vector<long long>> dp;            // dp[v][0 .. min(sz[v],V)]

void knap_dfs(int v, int p) {
    sz[v] = 1;
    dp[v].assign(V + 1, NEG_INF);        // NEG_INF = unreachable budget
    dp[v][0] = 0;                        // spend 0: seed with v's own term
    for (int c : adj[v]) {
        if (c == p) continue;            // parent guard = rooting
        knap_dfs(c, v);                  // child finalised first (post-order)
        int lim  = min(sz[v], V);        // budget reachable at v so far
        int clim = min(sz[c], V);        // budget child c can supply
        vector<long long> ndp(V + 1, NEG_INF);
        for (int i = 0; i <= lim; i++) if (dp[v][i] > NEG_INF)
            for (int j = 0; j <= clim && i + j <= V; j++)   // i + j <= V !!
                ndp[i + j] = max(ndp[i + j], dp[v][i] + dp[c][j]);
        dp[v].swap(ndp);
        sz[v] += sz[c];                  // running size grows by child's
    }
}
// after the call: answer is max_{j<=V} dp[root][j]
// counting version: replace max by + and NEG_INF seed by 0, product by sum.
\end{lstlisting}

\begin{complexitybox}
Time $\bigO{nV}$ by \thmref{c27-amortize}: the outer loop over children plus
the two capped inner loops is exactly the charged iteration count, $\bigO{nV}$,
and the post-order DFS touches each node once. Memory $\bigO{nV}$ for the state
arrays (or $\bigO{hV}$ if you free a child's array once merged). The stack
depth is $h$; at $n \ge 2\cdot10^{5}$ with chains legal, swap in the explicit
two-phase stack (Ch.12's \coderef{c09-iter}) so a path input does not overflow.
Verdict: the merge is cheap \emph{only because of the cap} -- forgetting
$\min(\cdot,V)$ turns a linear DP into a quadratic one.
\end{complexitybox}

\subsection{Worked Example: a $V=3$ Knapsack on $T$}

\begin{example}{Top-$j$ label-sum knapsack on the running tree $T$}{c27-knap-run}
Reuse $T$ from \figref{c27-runtree-knap}: node $v$ is an item of cost~$1$ and
value equal to its label; $\mathit{dp}[v][j]$ is the max value of a subset of
$\subtree(v)$ of size $\le j$. Since every item costs $1$, the answer for
$\mathit{dp}[v][j]$ is just the \emph{sum of the $j$ largest labels in
$\subtree(v)$} -- a closed form we can hand-check -- while the \emph{loop count}
of the merge is what \thmref{c27-amortize} is really about. Budget $V=3$. The
table below is computed by hand and reused, unchanged, by the trace box.
\end{example}

\begin{center}
\small
\renewcommand{\arraystretch}{1.2}
\begin{tabular}{l|cccccccccc}
\toprule
node $v$ & 1 & 2 & 3 & 4 & 5 & 6 & 7 & 8 & 9 & 10 \\
\midrule
$\size(v)$ & 10 & 4 & 5 & 1 & 2 & 1 & 3 & 1 & 2 & 1 \\
$\mathit{dp}[v][1]$ & 10 & 8 & 10 & 4 & 8 & 6 & 10 & 8 & 10 & 10 \\
$\mathit{dp}[v][2]$ & 19 & 13 & 19 & -- & 13 & -- & 19 & -- & 19 & -- \\
$\mathit{dp}[v][3]$ & 27 & 17 & 26 & -- & -- & -- & 26 & -- & -- & -- \\
\bottomrule
\end{tabular}
\end{center}

\noindent Sanity: $\mathit{dp}[1][3]=27 = 10{+}9{+}8$, the three largest labels in
all of $T$; $\mathit{dp}[3][3]=26=10{+}9{+}7$ (subtree of $3$ is
$\{3,6,7,9,10\}$, top three $=10,9,7$); a leaf has only $\mathit{dp}[\ell][1]$
(its subtree is one node), and node~$5$ (subtree $\{5,8\}$) only reaches $[\cdot][2]$.

\begin{figure}[H]
\centering
\begin{tikzpicture}[diagram, scale=0.82, transform shape,
    tnode/.style={circle, draw=primaryblue, line width=0.9pt, fill=primaryblue!8,
                  minimum size=7mm, font=\small\bfseries, text=primaryblue!60!black}]
  \node[tnode] (n1) at (0,0) {$1$};
  \node[tnode] (n2) at (-3,-1.6) {$2$};
  \node[tnode] (n3) at ( 3,-1.6) {$3$};
  \node[tnode] (n4) at (-5.2,-3.2) {$4$};
  \node[tnode] (n5) at (-1.6,-3.2) {$5$};
  \node[tnode] (n6) at ( 1.6,-3.2) {$6$};
  \node[tnode] (n7) at ( 5.2,-3.2) {$7$};
  \node[tnode] (n8) at (-1.6,-4.8) {$8$};
  \node[tnode] (n9) at ( 5.2,-4.8) {$9$};
  \node[tnode] (n10) at ( 5.2,-6.4) {$10$};
  \draw[tedge] (n1)--(n2) (n1)--(n3) (n2)--(n4) (n2)--(n5)
               (n3)--(n6) (n3)--(n7) (n5)--(n8) (n7)--(n9) (n9)--(n10);
  % merge cost labels on each parent->child edge: min(so far,V)*min(sz c,V)
  \node[annot, font=\scriptsize, text=warningred!75!black] at (-1.8,-0.6) {$1\!\cdot\!3{=}3$};
  \node[annot, font=\scriptsize, text=warningred!75!black] at ( 1.8,-0.6) {$3\!\cdot\!3{=}9$};
  \node[annot, font=\scriptsize, text=warningred!75!black] at (-4.4,-2.2) {$1\!\cdot\!1{=}1$};
  \node[annot, font=\scriptsize, text=warningred!75!black] at (-2.9,-2.2) {$2\!\cdot\!2{=}4$};
  \node[annot, font=\scriptsize, text=warningred!75!black] at ( 0.5,-2.2) {$1\!\cdot\!1{=}1$};
  \node[annot, font=\scriptsize, text=warningred!75!black] at ( 2.9,-2.2) {$2\!\cdot\!3{=}6$};
  \node[annot, font=\scriptsize, text=warningred!75!black] at (-1.2,-3.8) {$1\!\cdot\!1{=}1$};
  \node[annot, font=\scriptsize, text=warningred!75!black] at ( 4.6,-4.0) {$1\!\cdot\!2{=}2$};
  \node[annot, font=\scriptsize, text=warningred!75!black] at ( 4.6,-5.6) {$1\!\cdot\!1{=}1$};
  \node[annot, font=\scriptsize, text=difforange!75!black, anchor=west] at (5.9,-0.3)
    {edge $v\!\to\!c$:};
  \node[annot, font=\scriptsize, text=difforange!75!black, anchor=west] at (5.9,-1.0)
    {$\min(\text{so far},V)\cdot\min(\size(c),V)$};
  \node[annot, font=\scriptsize, text=successgreen!50!black, anchor=west] at (5.9,-2.1)
    {total $=3{+}9{+}1{+}4{+}1{+}6{+}1{+}2{+}1 = 28$};
  \node[annot, font=\scriptsize, text=successgreen!50!black, anchor=west] at (5.9,-2.8)
    {$\le Vn = 3\!\cdot\!10 = 30$};
\end{tikzpicture}
\caption{The \emph{merge cost ledger} for $T$ at $V=3$. Each non-root edge
carries the capped count of folding its child into the parent: the product of
$\min(S_{v,c},3)$ and $\min(\size(c),3)$. Over the nine edges the total is
$28 \le 30 = Vn$, exactly \thmref{c27-amortize}'s promise. Drop the cap (loop
heads at raw sizes) and the merge costs $45$ on $T$ -- and $\Theta(n^2)$ on a
chain.}
\label{fig:c27-mergetree}
\end{figure}

\begin{tracebox}[Merge ledger on $T$, $V=3$, children folded in numeric order]
post-order: 4,8,5,2,10,9,7,6,3,1. Leaves (4,6,8,10): zero children, cost $0$.\\
5: one child $8$, so-far$=1$ -> $\min(1,3)\!\cdot\!\min(1,3)=1$.\\
2: children $4$ then $5$. fold~4: $1\!\cdot\!\min(1,3)=1$; so-far$=2$.
   fold~5: $\min(2,3)\!\cdot\!\min(2,3)=4$. subtotal $5$.\\
9: child $10$, $1\!\cdot\!1=1$. 7: child $9$, $1\!\cdot\!\min(2,3)=2$.\\
3: children $6$ then $7$. fold~6: $1\!\cdot\!1=1$; so-far$=2$.
   fold~7: $\min(2,3)\!\cdot\!\min(3,3)=6$. subtotal $7$.\\
1: children $2$ then $3$. fold~2: $1\!\cdot\!\min(4,3)=3$; so-far$=5$.
   fold~3: $\min(5,3)\!\cdot\!\min(5,3)=9$. subtotal $12$.\\
Ledger: $1{+}5{+}1{+}2{+}7{+}12 = 28 \le 30$ \checkmark. Note node~1's first fold
costs only $3$ (so-far $=1$) though $\size(2)=4$ -- the running-size cap, in
action: the merge bound \thmref{c27-amortize} is $\bigO{nV}$, not $\bigO{n^2}$. \checkmark
\end{tracebox}

\begin{insightbox}[Key takeaways]
\begin{enumerate}[nosep, leftmargin=1.6em]
  \item \emph{Order is free, capping is mandatory.} Merging $2$ before $3$ at
        the root changes the per-edge numbers but never the $\bigO{nV}$ total.
        What must never change is the loop head $\min(\text{so far}, V)$.
  \item \emph{The result here is trivial to check.} Because costs are all $1$,
        $\mathit{dp}[v][j]$ is a top-$j$ sum -- you can verify the value in your
        head while the \emph{loop ledger} is the real lesson of 14.1.
  \item \emph{Real problems index by the count, not the value.} CF~855C below
        budgets ``how many type-$k$ vaults,'' a quantity $\le x \le 10$, and the
        state carries that count plus $v$'s own class -- exactly \defref{c27-knap-state}
        with a class dimension added.
\end{enumerate}
\end{insightbox}

\subsection{Variations and Twists}

\begin{itemize}[nosep]
  \item \textbf{Budget is ``number of special nodes'' (counting):} CF~855C and
        Summit CF~735E both spend the budget on \emph{chosen} vertices, so
        $j\in[0,V]$ is a count and the merge is $\sum\!\prod$ over splits --
        not $\max\!\sum$. The cap analysis is unchanged. \hfill\lv{3}
  \item \textbf{Exactly $j$, not at most $j$}: some problems want the budget
        spent to be exactly $j$; drop the $\le$ and keep unreachable slots at
        $\mathit{dp}=0$ (counting) or $-\infty$ (optimisation). Reading the
        answer becomes ``$\mathit{dp}[r][j]$ for the specific $j$,'' not a max. \hfill\lv{3}
  \item \textbf{Two coupled indices}: the state $\mathit{dp}[v][j][\text{cls}]$
        of CF~855C adds a constant ($\le 3$) class factor; it multiplies the
        merge constant, not the asymptotic -- still $\bigO{nV}$ with $V=x$. \hfill\lv{4}
  \item \textbf{When the budget is not small}: if $V$ ranges with the data
        (a value cap, a length, a time) the $\bigO{nV}$ may exceed a plain
        $\bigO{n\cdot\text{const}}$; \secref{c27-growing} designs the profile,
        and Ch.~24's centroid merge (\chref{c18}) takes over when $V \sim n$. \hfill\lv{4}
\end{itemize}

\subsection{Common Pitfalls}

\begin{warningbox}[Trap 1 -- uncapped loop head, quadratic in disguise]
The single most damaging budgeted-DP bug: writing the inner loop as
\lstinline{for (j = 0; j <= sz[c]; j++)} instead of
\lstinline{for (j = 0; j <= min(sz[c], V) && i + j <= V; j++)}. On the tiny
samples the answer is right; on a $2\cdot10^{5}$ chain the run time is
$\Theta(n^2)$ and you get a \tle. \thmref{c27-amortize} is not decoration --
the $\min(\cdot,V)$ cap is the hypothesis of the theorem.
\end{warningbox}

\begin{warningbox}[Trap 2 -- resetting the running size too late]
After folding a child, the ``size so far'' at $v$ must grow by that child's
$\size(c)$ \emph{before} the next child's $\min(S_{v,c},V)$ is evaluated.
Freeze $S$ at its initial value and every merge looks like $1 \cdot \min(\size(c),V)$;
you will underestimate the ledger on paper yet, worse, some codes seed
$\mathit{dp}[v]$ from the wrong ``reachable budget'' and drop valid splits --
silently wrong answers on wide nodes.
\end{warningbox}

\begin{warningbox}[Trap 3 -- overwriting \lstinline{dp[v]} in place]
The merge reads $\mathit{dp}[v][i]$ and writes $\mathit{dp}[v][i+j]$. If you
write into the same array you are folding a child into a partially-updated
parent, and a budget slot gets spent twice (the child's value counted with
itself). Always merge into a fresh $\mathit{dp}_{\text{new}}[v]$ for the child
and swap after -- as \lstinline{ndp} does in \coderef{c27-knap}.
\end{warningbox}

\subsection{Problem Lab}

\solvedproblem{c27-knap-solved}{Helga Hufflepuff's Cup}{Codeforces 855C}
  {https://codeforces.com/problemset/problem/855/C}{CF 1900 · L4}
  {budgeted subtree DP, three-class state, count of a special type as budget,
   modular counting, amortised O(nx) child-merge}

\begin{problemstatement}[Helga Hufflepuff's Cup (statement and samples verified against the live problem)]
The vault of Gringotts is a tree of $n$ rooms; room $v$ is assigned a
\emph{type}, an integer in $[1, m]$. Rooms of type $k$ have the \emph{highest}
security. A typing is valid if (i)~at most $x$ rooms have type $k$; and
(ii)~whenever a room has type $k$, every room \emph{adjacent} to it has type
strictly less than $k$. Count the valid typings of all $n$ rooms, modulo
$10^9+7$. Constraints: $1 \le n \le 10^5$, $1 \le m \le 10^9$,
$1 \le k \le m$, $1 \le x \le 10$. \textbf{Example~1}: $n=4, m=2$; edges
$1\text{--}2, 2\text{--}3, 1\text{--}4$; $k=1, x=2$; output \texttt{1} (type $1$
is highest, so a type-$1$ room forces neighbours of type $<1$ -- impossible --
hence no room is type $1$: every room is type $2$, exactly one way).
\textbf{Example~2}: $n=3, m=3$; edges $1\text{--}2, 1\text{--}3$; $k=2, x=1$;
output \texttt{13}. \textbf{Example~3}: $n=3, m=1$; edges $1\text{--}2, 1\text{--}3$;
$k=1, x=1$; output \texttt{0}.
\end{problemstatement}

\begin{attemptgate}[Attempt it yourself first -- \lv{4}, give yourself 35 minutes]
Two facts before you code. (a)~Only the \emph{relation to $k$} of a room's
type matters for the constraints -- collapse $m$ colours into three classes.
(b)~The budget $x$ is $\le 10$, so this is a $V=x$ knapsack whose item is
``this room is type $k$.''
\peek{What three classes do the $m$ types fall into relative to $k$, and how
      many actual types pick from in each class?}
\nudge{Below $k$: $k-1$ choices. Equal to $k$: exactly $1$. Above $k$: $m-k$.
      Call them L, C, H. Now: if a room is C, what \emph{must} its neighbours be?}
\fullhint{Class of a room is L (type $<k$, $k{-}1$ ways), C (type $=k$, $1$
way), H (type $>k$, $m{-}k$ ways). The adjacency rule becomes: a C room forces
\emph{all} neighbours to L; an H room may not touch a C room (so its neighbours
are L or H); an L room is free. Root the tree. Then state
$\mathit{dp}[v][j][\text{cls}]$ = number of valid typings of $\subtree(v)$ with
exactly $j$ type-$k$ rooms and $v$'s class $=$ cls. Merge children one at a
time over the split $j = i + j_c$ with $j \le x$ (\thmref{c27-amortize}, $V=x$):
if $v$ is C, each child must be L; if $v$ is L, a child is any of L,C,H; if $v$
is H, a child is L or H. Multiply by $\mathit{ch}[\text{cls}]$ at the leaf seed.
Answer $= \sum_{j\le x}\sum_{\text{cls}} \mathit{dp}[1][j][\text{cls}]$.}
\end{attemptgate}

\solstep{Understanding the Problem}
Three separate constraints are tangled together, and the trick is that they
\emph{decouple} once you look at classes instead of colours. The global
``at most $x$ rooms of type $k$'' is a \emph{budget} -- a count to spend -- so
it belongs on a DP index, and $x \le 10$ makes that index tiny. The local
``a type-$k$ room forces its neighbours below $k$'' is a \emph{parent--child
compatibility} between \emph{classes}, not colours: it only depends on whether
each end of an edge is L, C, or H. Once each room is reduced to one of three
classes, the number of original colourings is just the product of the per-class
option counts ($k{-}1$, $1$, $m{-}k$), and the combinatorics collapses into the
seed value at each node. What remains is a budgeted tree DP: the budget is the
count of C rooms.

\solstep{Key Observations}
\begin{observation}
A room's class is fully determined by $\operatorname{sign}(\text{type} - k)$,
and the adjacency rule in class language is: \textbf{C $\Rightarrow$ all
neighbours L}, \textbf{H $\Rightarrow$ no neighbour is C}, \textbf{L $\Rightarrow$
free}. These are edge-local and depend only on the two endpoints' classes -- the
rooted-DP compatibility the merge needs.
\end{observation}
\begin{observation}
In the rooted view the C/H rules only constrain the \emph{child} classes given
the parent class, so the constraint is a per-merge filter on the child's class
index: parent C admits only L children; parent H admits L and H; parent L
admits all. The reverse direction (child C forces parent L) is automatically
enforced because a C can only be placed under an L parent, and an L parent
admits C children -- the two readings agree.
\end{observation}
\begin{observation}
The count of C rooms is additive over subtrees and bounded by $x$: this is
exactly the budget index $j\in[0,x]$ of \defref{c27-knap-state}, so the merge is
a $\bigO{nx}$ convolution by \thmref{c27-amortize} with $V = x$.
\end{observation}

\solstep{Think It Through With Me}
\begin{thinkbox}
\deadend{Do a DP over the actual colours $1..m$; $m \le 10^9$ is fine with a
map or coordinate compression.}
{The colour value is irrelevant except relative to $k$; only three classes
matter. A colour-level state is both wrong (the adjacency rule speaks of
$<k$, $=k$, $>k$, not specific values) and astronomically expensive. Compress
\emph{first} -- that is the whole insight.}
\deadend{Fix the at-most-$x$ rooms of type $k$ by choosing which $\le x$ vertices
are C, then greedily colour the rest; multiply by $2^n$ for the two non-$k$ sides.}
{``Greedy colour the rest'' ignores that an H vertex touching a C vertex is
illegal, and the C set itself must be chosen so its neighbourhood is all L --
a \emph{structural} constraint that couples C-placement to the whole tree.
Also $2^n$ is fantasy: L has $k{-}1$ choices, H has $m{-}k$, they differ.}
\aha{Three classes, one budget index $j \le x$ for ``C rooms used,'' state
$\mathit{dp}[v][j][\text{cls}]$, children merged with the capped convolution
of \coderef{c27-knap}. Answer is the sum over $j\le x$ and cls at the root.}
\end{thinkbox}

\solstep{Solution}
Root at $1$. Precompute $\mathit{ch}[0]=k-1$ (L), $\mathit{ch}[1]=1$ (C),
$\mathit{ch}[2]=m-k$ (H), modulo $10^9+7$. At node $v$, seed
$\mathit{dp}[v][0][0]=\mathit{ch}[0]$, $\mathit{dp}[v][0][2]=\mathit{ch}[2]$,
and (if $x\ge1$) $\mathit{dp}[v][1][1]=\mathit{ch}[1]$. Then fold each child
$c$ with running budget $\le x$: for parent class C only child class L
contributes; for L all three; for H only child L and child H. After all children,
$\mathit{dp}[v][j][\text{cls}]$ counts valid subtree typings with exactly $j$
C rooms. Return $\sum_{j=0}^{x}\sum_{\text{cls}} \mathit{dp}[1][j][\text{cls}]$.

\solstep{Proof of Correctness}
\begin{proofbox}[\coderef{c27-855c-soln} counts exactly the valid typings]
By strong induction on the height of $\subtree(v)$ we show $\mathit{dp}[v][j][\text{cls}]$
equals the number of typings of $\subtree(v)$ that (a)~use exactly $j$ type-$k$
rooms, (b)~place $v$ in class $\text{cls}$, and (c)~satisfy the adjacency rule on
every edge \emph{internal} to $\subtree(v)$. \emph{Base (leaf)}: a single room in
class L has $k-1$ typings ($\mathit{ch}[0]$), $0$ C-rooms; in C it has $1$
type ($=k$) and $1$ C-room; in H it has $m-k$ types and $0$ C-rooms -- the three
seeds. No internal edge exists, so (c) is vacuous. \emph{Step}: assume the
claim for every child of $v$. A subtree typing of $v$ is $v$'s own class plus an
independent valid typing of each child subtree, and the C-count is additive:
$j = [\text{$v$ is C}] + \sum_c j_c$. The merge enumerates exactly these
splits $j = i + \sum_c j_c$ (children folded one at a time, budget $\le x$). The
only coupling between $v$ and a child $c$ is the adjacency rule on their shared
edge, which restricts $c$'s class to (L) when $v$=C, to (L or H) when $v$=H, and
to (L, C, or H) when $v$=L -- the three cases in the code. Hence the product over
children, summed over allowed class-splits, counts precisely the internal-rule
valid typings with $j$ C-rooms. The parent edge of $v$ is \emph{not} constrained
inside $\subtree(v)$; it is enforced higher up when $v$ is folded into its own
parent, which is why the state must keep $v$'s class. \emph{Answer}: at the root
there is no parent edge, so every internal edge is constrained, and summing over
all $j\le x$ and all classes counts exactly the typings obeying (i) ``$\le x$
type-$k$'' and (ii) the adjacency rule. \qed
\end{proofbox}

\solstep{Implementation}
\begin{lstlisting}[caption={Codeforces 855C -- three-class, C-count budget, one capped child-merge (\coderef{c27-knap} instantiated with a class dimension)}, label={lst:c27-855c-soln}]
#include <bits/stdc++.h>
using namespace std;
typedef long long ll;
const ll MOD = 1000000007LL;
// class: 0 = L (type < k), 1 = C (type == k), 2 = H (type > k)
int n, X; ll m, k; ll ch[3];
vector<vector<int>> adj;
vector<int> sz;
vector<vector<array<ll,3>>> dp;   // dp[v][j][cls], j = 0..X (C-rooms used)

void dfs(int v, int p) {
    sz[v] = 1;
    for (int j = 0; j <= X; j++) for (int c = 0; c < 3; c++) dp[v][j][c] = 0;
    dp[v][0][0] = ch[0];                 // v is L, contributes 0 C-rooms
    dp[v][0][2] = ch[2];                 // v is H, contributes 0 C-rooms
    if (X >= 1) dp[v][1][1] = ch[1];     // v is C, contributes 1 C-room
    for (int u : adj[v]) if (u != p) {
        dfs(u, v);
        vector<array<ll,3>> ndp(X + 1);
        for (int j = 0; j <= X; j++) for (int c = 0; c < 3; c++) ndp[j][c] = 0;
        for (int i = 0; i <= min(sz[v], X); i++)
            for (int ju = 0; ju <= min(sz[u], X) && i + ju <= X; ju++) {
                // v is L: child may be L, C, H
                for (int cu = 0; cu < 3; cu++)
                    ndp[i+ju][0] = (ndp[i+ju][0] + dp[v][i][0]*dp[u][ju][cu]) % MOD;
                // v is C: child forced L
                ndp[i+ju][1] = (ndp[i+ju][1] + dp[v][i][1]*dp[u][ju][0]) % MOD;
                // v is H: child may be L or H (not C)
                ndp[i+ju][2] = (ndp[i+ju][2] + dp[v][i][2]*dp[u][ju][0]) % MOD;
                ndp[i+ju][2] = (ndp[i+ju][2] + dp[v][i][2]*dp[u][ju][2]) % MOD;
            }
        dp[v] = ndp;
        sz[v] += sz[u];
    }
}
int main() {
    if (scanf("%d %lld", &n, &m) != 2) return 0;
    adj.assign(n + 1, {});
    for (int i = 0; i < n - 1; i++) {
        int u, v; scanf("%d %d", &u, &v);
        adj[u].push_back(v); adj[v].push_back(u);
    }
    scanf("%lld %d", &k, &X);
    ch[0] = (k - 1) % MOD; ch[1] = 1 % MOD; ch[2] = (m - k) % MOD;
    dp.assign(n + 1, vector<array<ll,3>>(X + 1));
    dfs(1, 0);
    ll ans = 0;
    for (int j = 0; j <= X; j++) for (int c = 0; c < 3; c++) ans = (ans + dp[1][j][c]) % MOD;
    printf("%lld\n", ans);
    return 0;
}
\end{lstlisting}

\solstep{Trace}
\begin{tracebox}[Official Example~2: $n=3, m=3$, edges $1\text{--}2, 1\text{--}3$, $k=2, x=1$ $\Rightarrow$ 13]
Here $\mathit{ch}[0]=k{-}1=1$, $\mathit{ch}[1]=1$, $\mathit{ch}[2]=m{-}k=1$: each
class has exactly one colour, so counting class-assignments $=$ counting typings.
Root $1$, leaves $2,3$. Each leaf $u$: $\mathit{dp}[u][0][0]=1$, $\mathit{dp}[u][0][2]=1$,
$\mathit{dp}[u][1][1]=1$, all else $0$.\;
Node $1$ seed: $[0]\!\to\!(\text{L}1,\text{C}0,\text{H}1)$, $[1]\!\to\!(\text{L}0,\text{C}1,\text{H}0)$.\ 
Budget cap $i+j\le x=1$, so only $(i,ju)\in\{(0,0),(0,1),(1,0)\}$.\\[6pt]
\textbf{Fold leaf~2.} $(0,0)$: L$\to1\!\cdot\!(1{+}0{+}1)=2$, H$\to1\!\cdot\!1{+}1\!\cdot\!1=2$;
$(0,1)$: L$\to1\!\cdot\!1=1$ (child C); $(1,0)$: C$\to1\!\cdot\!1=1$ (parent C, child L).
State after: $[0]\!\to\!(2,0,2)$, $[1]\!\to\!(1,1,0)$.\\[6pt]
\textbf{Fold leaf~3.} $(0,0)$: L$\to2\!\cdot\!2=4$, H$\to2\!\cdot\!1{+}2\!\cdot\!1=4$;
$(0,1)$: L$\to2\!\cdot\!1=2$; $(1,0)$: L$\to1\!\cdot\!2=2$, C$\to1\!\cdot\!1=1$.
Final: $[0]\!\to\!(\text{L}4,\text{C}0,\text{H}4)$, $[1]\!\to\!(\text{L}4,\text{C}1,\text{H}0)$.\\[6pt]
\textbf{Answer} $=\sum_{j\le1}\sum_{\text{cls}}$: $j{=}0$ gives $4{+}0{+}4=\mathbf{8}$
(zero C-rooms: every node L or H $=2^3$); $j{=}1$ gives $4{+}1{+}0=\mathbf{5}$
(exactly one C: root C forces both leaves L $=1$ way, or a leaf C whose parent is
L and whose sibling is free L/H $=2$ ways $\times2=4$). Total $8+5=\boxed{13}$
\checkmark\ matches CF Example~2.
\end{tracebox}

\begin{tracebox}[Examples 1 and 3 fold the same rule]
Ex.~1 ($m=2, k=1$): $\mathit{ch}[0]=k{-}1=0$, so any L-forced state is $0$. A C
room forces all neighbours to L ($0$ ways) $\Rightarrow$ no C is usable; every room
must be H ($\mathit{ch}[2]=m{-}k=1$) $\Rightarrow$ exactly one typing. Output
\texttt{1}. \checkmark\\
Ex.~3 ($m=1, k=1$): only type $1=C$ exists, $\mathit{ch}[0]=\mathit{ch}[2]=0$.
Placing a C forces its neighbours to L ($0$ ways) on every edge with $n\ge2$, so no
valid typing exists. Output \texttt{0}. \checkmark
\end{tracebox}

\solstep{Complexity}
\begin{complexitybox}
Time $\bigO{nx}$ by \thmref{c27-amortize} with $V = x \le 10$: each node's array
has length $\le x+1$, children merge with the $i + j \le x$ cap, and the constant
factor is the (parent class, child class) pairs ($\le 3\!\cdot\!3$). For $n=10^5$
that is $\le 30\!\cdot\!10^{6}$ operations. Memory $\bigO{nx}$ for the DP arrays
($\sim\!10^6$ \lstinline{ll}) plus $\bigO{n}$ adjacency. Everything is modulo
$10^9+7$ -- keep the accumulators in \lstinline{long long} and take $\%$ MOD inside
the inner loop, since a $k-1$ factor is up to $10^9$. Verdict: \ac; this is the
archetypal budgeted tree knapsack -- three classes, one count index, one capped
child-merge.
\end{complexitybox}

\solstep{Pitfalls and Variants}
\begin{warningbox}
Four ways to lose CF~855C: (i)~\emph{forget the modulo inside the inner loop} --
$\mathit{ch}[0]=k-1$ is up to $10^9$ and products overflow \lstinline{int} and even
double the modulus before you reduce; (ii)~\emph{let an H parent keep a C child}
-- the rule ``H has no C neighbour'' is exactly the class filter in the merge, drop
it and Example~2 over-counts; (iii)~\emph{read only $\mathit{dp}[1][x][\cdot]$}
instead of summing $j\le x$ -- the bound is ``\emph{at most} $x$''; (iv)~loop head
at $\size(v)$ not $\min(\size(v),x)$ -- correct answer, \tle on the big tests.
Variant: swap ``exactly $j$ C-rooms'' for a second index counting H-rooms and you
get the shape of the Summit CF~735E, where the budget is distance-to-nearest-black
instead of a global count -- same merge discipline.
\end{warningbox}

\subsection{Practice Problems}

\begin{stuckbox}
\begin{itemize}[nosep, leftmargin=1.4em]
  \item Did your state carry the \emph{class} of $v$ so the parent-edge rule can
        be checked when $v$ is folded into its parent? Without it the C-forces-L
        constraint is invisible across the subtree boundary.
  \item Are the loop heads $\min(\text{size so far}, x)$ and $\min(\size(c), x)$
        with $i + j \le x$? That cap is \thmref{c27-amortize}'s hypothesis.
  \item Did you sum the answer over \emph{all} $j \le x$ and all three classes at
        the root, and reduce mod $10^9+7$ on every product?
\end{itemize}
\end{stuckbox}

\begin{practiceset}
  \practice{Apple Tree}{Codeforces 461B}
    {https://codeforces.com/problemset/problem/461/B}{CF 2100 · L5}
    {Delete a set of edges so every resulting component holds \emph{exactly one}
    black vertex; count the sets mod $10^9+7$. This is \defref{c27-knap-state}
    with budget $V=1$: $\mathit{dp}[v][0/1]$ = ``the component still open at $v$
    already owns a black vertex / does not yet.'' Merging a child is a two-term
    knapsack (cut the edge $\to$ the child's part must already own its black;
    keep it $\to$ the black counts add, capped at~$1$). Write the merge with the
    same $\min(\text{so far},V)$ discipline as \coderef{c27-knap}. Samples
    (verified): $3$ vertices $\to 2$, $6\to1$, $10\to27$.
    }
\end{practiceset}

\paragraph{Back-cites (URLs live elsewhere in the book).}
\begin{itemize}[nosep,leftmargin=1.4em]
  \item \textbf{CF~161D ``Distance in Tree''} (\chref{c10} owns this URL): the
        distance-$k$ pair count is the histogram merge Ch.~13 used \emph{without}
        proving $\bigO{nk}$; \thmref{c27-amortize} is that missing proof, read
        with $V = k$. Re-derive CF~161D's $\mathit{dp}[v][d]$ array from
        \coderef{c27-knap} and the amortisation now justifies the $10^5 \cdot 500$
        budget.
  \item \textbf{CSES ``Tree Matching'' family / CF~570D} (\chref{c10} owns
        CF~570D): the depth-capped subtree query DP is the same budgeted state
        answering per-root; that is \secref{c27-knapsack}'s array read at every
        vertex, and Chapter~15 (\chref{c20}) is what makes the ``every root''
        reading cheap -- not this chapter.
\end{itemize}

\begin{recallbox}[10-second refresher -- 14.1 Tree knapsack and the merge bound]
The state is an array $\mathit{dp}[v][0..V]$ over a budget; children are merged
\emph{one at a time} by a convolution capped at $\min(\text{size so far},V)$ and
$\min(\size(c),V)$ with $i+j\le V$. \thmref{c27-amortize} charges each iteration
to a node-pair at their LCA -- each pair once, $\le V$ per node -- so the whole
DP is $\bigO{nV}$. CF~855C (\probref{c27-knap-solved}) is three classes (L/C/H
around $k$) and one count index $j\le x$; the ledger on $T$ (root $\mathit{dp}[1][3]=27$,
$28$ merge iterations $\le 3n$) is the sanity check.
\selftest{
  \item Why can Ch.~13's constant-tuple fold be $\bigO{n}$ but the array state
        needs a real amortisation argument?
  \item In \coderef{c27-knap}, name the \emph{two} places the cap $\min(\cdot,V)$
        appears and what each prevents.
  \item CF~855C: give the three class-option counts and the child-class filter
        for a parent of class H.
  \item State \thmref{c27-amortize}'s charge in one sentence: to what object is a
        merge iteration charged, and why is it charged at most once?
  \item Reading the answer as $\mathit{dp}[1][x][\cdot]$ instead of
        $\sum_{j\le x}$ is which trap?
}
\end{recallbox}
